% Paper 4, §1 — Introduction
% Thesis: interpretability is method-plural because the underlying object —
% a frontier transformer's internal computation — is plural. Five methods
% (lens, probes, SAEs, activation patching, circuit tracing) commit
% differently along four axes: correlational vs causal, component- vs
% feature-level, supervised vs unsupervised, static-readout vs intervention.
% "Partially read" claim: human-legible structure for some behaviors, on some
% models, with cross-method agreement on some claims. Closing question for
% §11/§12: when do two methods *have* to agree, and what does it mean when
% they don't?

\section{Introduction}
\label{sec:introduction}

Interpretability is method-plural. A modern frontier transformer is, at
the level of weights and activations, a stack of attention and MLP
sublayers writing into a residual stream of dimension $D \in \{2048,
4096, 8192, 16384\}$ over sequences of length $S$ up to $\sim10^6$
tokens, with parameter counts from $10^9$ to $10^{12}$ and a vocabulary
of $|\mathcal{V}| \in \{32\text{k}, 128\text{k}, 256\text{k}\}$ tokens
(Paper~1 \S2). The mechanistic-interpretability program asks what that
stack is \emph{computing} --- not just what it predicts, but what
intermediate representations it builds, which of those representations
it uses, and how those representations compose into the behaviors we
care about. The field has not converged on a single method for
answering that question, and this paper argues it should not. The
underlying object is plural: the same residual-stream activation
$\mathbf{h}^{(\ell, t)} \in \mathbb{R}^D$ at layer $\ell$, position
$t$ is simultaneously a partial vocabulary distribution
\citep{belrose2023-tuned-lens}, a vector of linearly decodable
properties \citep{marks-tegmark-2023-truth}, a sparse combination of
overcomplete features \citep{gao2024-topk-saes}, a node in a causal
graph \citep{meng2022-rome}, and a participant in a sparse
behavior-implementing subgraph \citep{wang2022-ioi}. Each of those is
a real fact about $\mathbf{h}^{(\ell, t)}$, observable with a
different instrument, and load-bearing for a different question. The
methodological pluralism is principled, not transitional.

This paper is a method-by-method comparison of the five interpretability
techniques that have stabilized into the field's canon between roughly
2017 and 2026. Method-pluralism is the throughline, but the
load-bearing claim is sharper: the underlying transformer's internal
computation can be \emph{partially} read --- we can extract
human-legible structure for \emph{some} behaviors, on \emph{some}
models, with cross-method agreement on \emph{some} claims --- and the
field is honest about which behaviors, which models, and which claims.
Complete reverse-engineering of a frontier model is not on the table
in 2026; partial, locally-faithful, cross-method-validated explanation
is.

\subsection{Five methods, four commitments}
\label{sec:five-methods}

Five interpretability methods anchor the literature, in roughly the
order they crystallized:

\begin{description}
  \item[Lens techniques (\S\ref{sec:logit-lens-family}).] Project an
    intermediate residual-stream state $\mathbf{h}^{(\ell, t)}$ through
    the model's own pre-unembed normalization and unembedding $W_U$ to
    obtain a layer-$\ell$ ``as-if-final'' distribution over
    $\mathcal{V}$. The \emph{logit lens}
    \citep{nostalgebraist2020-logit-lens} is the no-training baseline;
    the \emph{tuned lens} \citep{belrose2023-tuned-lens} fixes
    documented cross-model failures with a per-layer affine translator.
  \item[Probes (\S\ref{sec:probing}).] Train a (typically linear)
    classifier $\hat y = \sigma(\mathbf{w}^\top \mathbf{h}^{(\ell, t)} +
    b)$ on \emph{frozen} activations to predict an external property
    $y$ \citep{alain-bengio-2017}. High accuracy is the (correlational)
    signal that $y$ is linearly decodable from $\mathbf{h}^{(\ell, t)}$.
  \item[Sparse autoencoders (\S\ref{sec:sparse-autoencoders}).]
    Decompose activations $\mathbf{x} \in \mathbb{R}^n$ into a sparse
    linear combination of $M \gg n$ learned dictionary directions
    $\{\mathbf{d}_i\}$ via an encoder--decoder pair under a joint
    reconstruction-plus-sparsity loss \citep{gao2024-topk-saes}. The
    commitment: features are linear, additive, and identifiable from
    activations alone, without supervision.
  \item[Activation patching (\S\ref{sec:patching}).] The canonical
    causal intervention. The three-run protocol of
    \citet{meng2022-rome} caches the clean run, records the corrupted
    run, then substitutes one component's clean activation into the
    corrupted run; the indirect effect $\mathrm{IE}_C = m_{\text{pt}} -
    m_{*}$ is the causal estimator \citep{zhang2023-apatching}. ROME
    \citep{meng2022-rome} extends this from diagnostic to weight-level
    surgery (\S\ref{sec:rome}).
  \item[Circuit tracing (\S\ref{sec:circuits}).] Compose the previous
    methods into a sparse, causal subgraph $C \subseteq M$ of the
    model's full computational graph that explains an identifiable
    capability. \citet{wang2022-ioi} hand-reverse-engineered
    indirect-object identification in GPT-2 small;
    \citet{conmy2023-acdc} automated the discovery as iterative
    edge-pruning; \citet{lindsey2025-circuit-tracing} scaled the
    program to production Claude 3.5 Haiku via cross-layer transcoders.
\end{description}

These five do not partition cleanly along a single axis. They commit
along (at least) four:

\begin{enumerate}
  \item \textbf{Correlational vs causal.} Lenses and probes read off
    activations and report what is \emph{decodable}; activation
    patching, ROME, and circuit tracing intervene on activations or
    weights and report what is \emph{used}. The correlation--causation
    gap (\S\ref{sec:correlation-causation}) is the
    methodological seam: a high-accuracy probe is consistent with the
    decoded property being computed and consumed, computed and
    discarded (\emph{vestigial information}), or being a confound
    outside the model entirely \citep{belinkov2022-probing-survey}.
  \item \textbf{Component-level vs feature-level.} Activation patching
    and circuit-tracing-via-attention-heads identify
    \emph{components} (heads, MLPs, hidden states) that mediate a
    behavior; SAEs and feature-level circuit tracing identify
    \emph{directions} in residual-stream space, which can be smaller
    or larger than any component basis. The two granularities answer
    different questions.
  \item \textbf{Supervised vs unsupervised.} Probes train against
    labels; SAEs train against the activations themselves. Lenses
    train against the model's own final-layer output (a form of
    self-supervision). The supervision signal determines which
    properties the method can find: a probe finds what it is asked to
    find; an SAE finds what its sparsity prior considers a feature.
  \item \textbf{Static readout vs intervention.} Lenses, probes, and
    SAEs are static --- they compute on cached activations without
    perturbing the forward pass. Activation patching, ROME, and
    circuit-validation knockouts are interventional. The static
    methods scale gracefully; the interventional methods carry the
    causal load.
\end{enumerate}

The methods are not strictly ranked. Lenses are the cheapest readout
($L \cdot (D^2 + D)$ tuned-lens parameters --- under $10^9$ for
Pythia~12B \citep{belrose2023-tuned-lens}); probes are next; SAEs sit
at training-pretraining-fraction compute (Gemma Scope at $\sim20\%$ of
GPT-3 pretraining \citep{gemma-scope-2024}); activation patching costs
$O(N_{\text{components}})$ forward passes per input pair; circuit
tracing on production models inherits the SAE/transcoder training cost
\emph{plus} the per-task graph-pruning cost. The compute envelope is
itself a methodological commitment: a method that is too expensive to
run does not enter the field's working vocabulary.

\begin{intuition}
The five methods are five different instruments aimed at the same
residual stream. The lens reads what the model is about to predict;
the probe reads what the model has computed; the SAE reads how the
model has decomposed what it computed; activation patching tests
which of those computations the model uses; circuit tracing
composes all four into a sparse, causal explanation of a specific
behavior. Returning to canonical form: each method is a different
projection of the activation tensor $\mathbf{h}^{(\ell, t)} \in
\R^{\BSH}$, and ``method-plural'' is the
recognition that no single projection is sufficient.
\end{intuition}

\subsection{What ``partially read'' means}
\label{sec:partially-read}

The thesis sentence --- \emph{the internal computation can be
partially read} --- has three quantifiers. \emph{Some behaviors}: the
GPT-2-small IOI circuit \citep{wang2022-ioi}, GPT-2-XL factual recall
\citep{meng2022-rome}, and selected Claude 3.5 Haiku capabilities
\citep{lindsey2025-circuit-tracing} are reverse-engineered to a
mechanism-level account; almost everything else a frontier model does
is not. \emph{Some models}: SAEs are fully released for the Gemma~2
family at 2B/9B/27B \citep{gemma-scope-2024}; tuned lenses cover the
GPT-2 / Pythia / LLaMA-1 lineage \citep{belrose2023-tuned-lens} but
have not all been published for Llama 3/4, Qwen 2.5/3, DeepSeek-V3, or
Gemma 3 (\S\ref{sec:scaling-snapshot}). \emph{Some claims}: the
mid-layer-MLP localization of factual recall, the IOI circuit's
26-head decomposition, and the linear decodability of truth on
LLaMA-family models are claims multiple methods agree on; the
canonicality of SAE features
\citep{leask2025-sae-not-canonical}, the existence of a single truth
direction across datasets \citep{marks-tegmark-2023-truth}, and the
reproducibility of attribution-graph findings outside the originating
lab are claims the methods disagree on.

The closing question of this paper, taken up in
\S\ref{sec:cross-method}, is the diagnostic version of that disagreement:
\emph{when do two methods have to agree, and what does it mean when
they don't?} When two methods answer the same question about the same
mechanism --- e.g., activation patching and ACDC on IOI --- they
\emph{have} to agree, and divergence is a falsification signal for at
least one method's assumptions. When two methods answer different
questions --- e.g., a probe (decodable?) and patching (used?) --- they
\emph{can} disagree without either being wrong, and the disagreement
is principled. Distinguishing the two cases is a load-bearing
methodological skill that this paper tries to teach.

\subsection{Roadmap}
\label{sec:roadmap}

The paper is organized to make the method-pluralism legible.
\S\ref{sec:substrate} fixes the substrate and notation: the residual
stream as a linear additive bus, attention as QK+OV, MLPs as candidate
key--value associative memories, the unembedding $W_U$ as the
canonical readout, and the formal definitions of \emph{feature} (a
direction $\mathbf{d} \in \mathbb{R}^D$) and \emph{circuit} (a
subgraph satisfying \citet{wang2022-ioi}'s
faithfulness/completeness/minimality criteria) used throughout.

\S\S\ref{sec:logit-lens-family}--\ref{sec:circuits} present the five
methods and their immediate critical literature.
\S\ref{sec:logit-lens-family} covers lenses; \S\ref{sec:probing}
covers probes; \S\ref{sec:correlation-causation} formalizes the
correlation--causation gap and the activation-intervention
resolution; \S\ref{sec:truth-directions} works the truth-direction
case study; \S\ref{sec:sparse-autoencoders} covers SAE architectures;
\S\ref{sec:sae-canonicality} takes up the canonicality question raised
by \citet{leask2025-sae-not-canonical}; \S\ref{sec:patching} covers
activation patching; \S\ref{sec:rome} covers ROME and the
interpretability-to-surgery bridge; \S\ref{sec:circuits} covers
circuit tracing from IOI to Haiku.

\S\S\ref{sec:cross-method}--\ref{sec:contradictions} are the
cross-method synthesis. \S\ref{sec:cross-method} catalogs the
convergence and divergence cases method-by-method and behavior-by-behavior;
\S\ref{sec:scaling-snapshot} surveys what scales (SAEs, probes,
attribution graphs in narrow regimes) and what does not yet (manual
circuit hunts, full activation patching at frontier scale, lens
coverage of MLA / MoE / NSA architectures), with a 2026 frontier
snapshot table of which method covers which model;
\S\ref{sec:contradictions} concentrates the live disagreements ---
SAE canonicality, lens cross-model unreliability, probe correlation
vs causation, truth-direction multiplicity, circuit-tracing
reproducibility past Anthropic-2025 --- and names, for each, the
experiment that would resolve it.

The reader who wants the headline of each method in isolation can
read \S\ref{sec:substrate} and any single
\S\ref{sec:logit-lens-family}--\S\ref{sec:circuits}. The reader who
wants the comparative argument should read
\S\ref{sec:substrate} through \S\ref{sec:contradictions} in order:
the cross-method evidence is what makes ``the internal computation
can be partially read'' a defensible claim rather than a hopeful
slogan, and that evidence is built incrementally across the paper's
arc.
